\documentclass[12pt, letterpaper]{article}
\usepackage[titletoc,title]{appendix}
\usepackage{color}
\usepackage{booktabs}
\usepackage[usenames,dvipsnames,svgnames,table]{xcolor}
\definecolor{dark-red}{rgb}{0.75,0.10,0.10} 
\usepackage[margin=1in]{geometry}
\usepackage{dcolumn}
\usepackage[linkcolor=dark-red,
colorlinks=true,
urlcolor=blue,
pdfstartview={XYZ null null 1.00},
pdfpagemode=UseNone,
citecolor={dark-red},
pdftitle={Coming to Dislike}]{hyperref}
%\usepackage{biblatex}
\usepackage{multibib}
\newcites{sec}{Other Sources}
\usepackage{geometry} % see geometry.pdf on how to lay out the page. There's lots.
\geometry{letterpaper}               % This is 8.5x11 paper. Options are a4paper or a5paper or other... 
\usepackage{graphicx}                % Handles inclusion of major graphics formats and allows use of 
\usepackage{amsfonts,amssymb,amsbsy}
\usepackage{amsxtra}
\usepackage{natbib}
\usepackage{verbatim}
\setcitestyle{round,semicolon,aysep={},yysep={;}}
\usepackage{setspace}            % Permits line spacing control. Options are \doublespacing, \onehalfspace
\usepackage{sectsty}             % Permits control of section header styles
\usepackage{lscape}
\usepackage{fancyhdr}            % Permits header customization. See header section below.
\usepackage{url}                     % Correctly formats URLs with the \url{} tag
\usepackage{fullpage}       %1-inch margins
\usepackage{multirow}
\usepackage{rotating}
\setlength{\parindent}{3em}

\usepackage[T1]{fontenc}
\usepackage{bm}
\usepackage{libertine}

\usepackage{chngcntr}

% Caption
\usepackage[hang, font=small,skip=0pt, labelfont={bf}]{caption}
\captionsetup[subtable]{font=small,skip=0pt}
\usepackage{subcaption}

% tt font issues
% \renewcommand*{\ttdefault}{qcr}
\renewcommand{\ttdefault}{pcr}

\usepackage{lscape}
\renewcommand{\textfraction}{0}
\renewcommand{\topfraction}{0.95}
\renewcommand{\bottomfraction}{0.95}
\renewcommand{\floatpagefraction}{0.40}
\setcounter{totalnumber}{5}
\makeatletter
\providecommand\phantomcaption{\caption@refstepcounter\@captype}
\makeatother

\title{Coming to Dislike Your Opponents:\\The Polarizing Impact of Political Campaigns}
\date{\vspace{-5ex}}
\author{Gaurav Sood\thanks{Gaurav can be reached at \href{mailto:gsood07@gmail.com}{\texttt{gsood07@gmail.com}}.} \and Shanto Iyengar\thanks{Shanto Iyengar is Chandler Chair in Communication and Professor of Political Science at Stanford University. He can be reached at: \href{mailto:siyengar@stanford.edu}{\texttt{siyengar@stanford.edu}}.}}


\begin{comment}
	 
setwd(paste0(basedir, "polarcamp/ms/"))
tools::texi2dvi("ComingToDislike.tex", pdf=TRUE,clean=TRUE)
setwd(basedir)

\end{comment}

%===================Abstract======================= 

\begin{document}
\maketitle
\doublespacing
%===================Abstract======================= 
\thispagestyle{empty}
\begin{abstract}
\noindent
Loud and vitriolic campaigns are increasingly the norm. For an electorate for which partisanship is a salient social identity, campaign messages questioning the intentions, integrity, and patriotism of political opponents are liable to not only reinforce partisans' stereotypes of the other side, but also engender new negative stereotypes.  We use data from multiple large national surveys, and the Wisconsin Advertising Project to demonstrate that partisans' evaluations of their opponents become more negative over the course of the campaign. Exposure to televised political advertising, especially negative advertising, increases partisan affect. We discuss the implications of our findings for current debates about the extent of partisan polarization within the mass public, and the consequences of such polarization for electoral accountability.
\end{abstract}

\newpage

A new chill has descended on inter-partisan relations. Republicans increasingly dislike Democrats, and think most Democrats to be ``hypocritical'', and ``closed-minded''; for their part, Democrats feel much the same way about Republicans (\citealt{iyengar2012}, see also \citealt{Iyengar2014, Mason2013}).  Not only are partisans increasingly willing to ascribe negative traits to supporters of the opposing party, they are also increasingly socially distant to them.  For instance, large proportions of both Democrats and Republicans are ``unhappy'' at the thought of their relation marrying a supporter of the main opposing party \citep{iyengar2012}. 

It is against this backdrop of strong affective partisan attachments that the increasingly long, loud, and hostile political campaigns are run. Campaigning in presidential elections now starts in earnest a good two years in advance of the election. On television alone, hundreds of thousands of ads are aired; over a million ads were aired as part of the 2012 general election campaign.\footnote{Wesleyan Media Project: \href{http://mediaproject.wesleyan.edu/2012/11/02/presidential-ad-war-tops-1m-airings/}{http://mediaproject.wesleyan.edu/2012/11/02/presidential-ad-war-tops-1m-airings/}} Candidates and their surrogates air primarily negative advertisements, and news reports obsess about the most vicious of the negative ads, giving the most perverse of the negative advertisements even greater shelf life \citep{geer2010}. For instance, during the period preceding the 2004 presidential election, there were more news reports about the notorious Swift Boat Ad questioning Senator Kerry's military record than on the war in Iraq \citep{geer2010}.

Partisans process this avalanche of political messages in a manner that favors the party they identify with \citep{chang2003, lodge1986, rahn1993}. As a result, partisans' already warm feelings for their party's candidate increase even more, while their dislike for the opposing party candidate takes an even sharper negative turn. Or so we will show. As part of our exploration into how political campaigns affect \emph{affect} towards candidates, we also show that voters learn little about the ideological positions of the candidates, precluding learning as the primary reason for polarization of candidate assessments over the course of the campaign.

We assess the role of campaigns as engines of polarization of candidate assessments using survey data from the 2000, 2004 and 2008 National Annenberg Election Studies (NAES), and televised presidential advertisement data for the years 2000 and 2004 from the Wisconsin Advertising Project. We test whether partisans' evaluations of their opponents become more hostile as a consequence of exposure to the campaign in general and to negative advertising in particular. The data suggest a strong connection between negative campaigns and polarization of feelings towards the candidates.

We begin with an overview of our theoretical expectations, derived from social identity theory and decades of research showing the reinforcing effects of political campaigns.  Next, we describe our data sources, research design, and measurement strategy. Following which, we present the results. Finally, we discuss the implications of our results for ongoing debates over the causes and consequences of mass polarization.  

\section*{Partisan Processing of Information}
Republicans increasingly dislike Democrats, and Democrats feel much the same way about Republicans \citep{iyengar2012}. One salient consequence of partisan affect is that it acts as ``perceptual screen'' \citep{campbell1960}, encouraging selective exposure and processing of political messages that people receive during campaigns. It is well documented that voters react to political messages not as dispassionate observers, but as biased partisans \citep{chang2003, rahn1993, lodge1986, campbell1960, lord1979}.  For instance, no matter how flawed the actual performance of candidates during televised debates, partisans are quick to declare ``their'' candidate the winner \citep{sigelman1984, iyengar2000b}. For example, a majority of the Republicans felt that President Ford had out-debated Jimmy Carter despite Ford's repeated gaffes concerning the autonomy of Eastern Europe in the second debate \citep{sears1979}.

In case of television ads, which are transparently one-sided, making it relatively easy for partisans to recognize partisanship of the sponsor and react accordingly, partisan filtering is starker still. Partisans accept or reject ads almost instantly based on sponsor of the ad \citep{iyengar2008, mcguire1969} for the classic discussion of acceptance-rejection factors in the persuasion process, see \citealt{mcguire1985}). In a series of experiments, Ansolabehere and Iyengar demonstrated that ads proved persuasive only among voters who shared the partisanship of the sponsoring candidate.  Exposure to a single advertisement boosted support for the sponsor by 14 percent among in-partisans, but by only 3 percent among independents and out-partisans \citep[][p. 76]{ansolabehere1995}.  In the case of in-partisans with little interest in politics, the sponsor's share of the vote actually increased by 25 percent (p. 80). Similar experimental findings emerged in studies of the 1992 and 1996 presidential campaigns which confirmed that the reinforcing effect of exposure to campaign ads was especially pronounced among younger voters, who are typically less partisan \citep{iyengar2000a}.  In short, as the scope and volume of ad campaigns increases, partisans of either side are more likely to line up in their respective corners \citep[for a review see][]{iyengar2000b}.

\section*{A Preference for Negative Advertising}
While exposure to any political advertising may be expected to increase partisan affect, we contend that negative ads are especially effective in increasing partisan affect.  A strong negativity bias influences information processing \citep[see][]{soroka2012} making people more likely to attend to negative rather than positive appeals \citep{meffert2006}.  Increased attention to negative messages makes them more memorable \citep{lang1991, bradley2007, chang2001, newhagen1991}. Moreover, negative ads are more effective at eliciting short term aversive sentiment against candidates than positive ads \citep{chang2001}.  Given these effects of differential attention, retention, and short-term affective response, we think it especially likely that negative ads will prove more effective than positive ads in evoking partisan affect.  This expectation is entirely consistent with the conventional wisdom among practitioners and journalists about efficacy of negative campaigns. There is also some research that suggests that negative ads lower evaluations of the target \citep{gaines2011}.

\section*{Learning}
Campaigns may also increase awareness about policy differences between the candidates and the parties. Since a greater share of negative ads than positive ads broaches issues---about two-thirds of the negative ads in some previous campaigns have been issue ads \citep[][p. 62]{geer2008}, \textemdash campaign season may be seen as an (ever larger) avalanche of information about issue positions. In principle, this avalanche of information should leave even the comatose informed. Yet the information that has surfaced until now is about modest learning of issue positions \citep{geer2008, lau2007, lau2009, zhao1995, ridout2004, ansolabehere1995}. ``[T]his effect is small, featuring unadjusted and adjusted mean effect sizes in the single digits; across the 15 pertinent findings, the average boost in campaign knowledge associated with negative campaigning amounts to less than one-tenth of a standard deviation'' \citep{lau2007}.

So we think it unlikely that it is increased awareness of issue positions that is the primary explanation for increased affect. Furthermore, should voters become informed about the policy positions of the candidates \textemdash typically more extreme than that of the public \citep{fiorina2012} \textemdash affection towards own party's candidate should attenuate, and disaffection with other party's candidate should grow even stronger. It is possible, though yet largely unexplored, that campaigns, perhaps by equal measure, increase awareness of incorrect information, and that it is increases in misinformation that are behind increased difference between affect towards in-party and out-party candidate. 

\section*{Data and Research Design}
Our empirical case rests upon the NAES Rolling Cross Section (RCS) data for the years 2000, 2004, and 2008, and television advertisement data from the Wisconsin Advertising Project for the years 2000 and 2004. Using the NAES RCS data, we track how partisan affect evolves over the course of the campaign among self-identified partisans. Next, to identify the role of the campaign more directly, we compare trends in affect in battleground vis-\`{a}-vis non-battleground states.  Presidential campaigns are geared towards winning the Electoral College. Given the winner-takes-all design of presidential vote at the state level, and differences in the extent to which each state contributes to the Electoral College, presidential campaign volume is much higher in electorally competitive states. As a consequence similar people are exposed to vastly different levels of campaigning depending on where they are located \citep[see][]{ashworth2006, ansolabehere2006}. Given greater `dosage' of the treatment in battleground states, any observed effects should be more pronounced there.

Differences between battleground and non-battleground states, adjusted for covariates, give us an estimate of the effect of total campaign activity across the two. To parcel out the impact of televised campaign advertising, we use an identification strategy used by \citet{huber2007}. The identification strategy rests upon the less than precise geographical correspondence between designated media market areas (DMAs) and state boundaries. A DMA in a battleground state sometimes encompasses parts of neighboring non-battleground state(s).  Residents of the latter can therefore be considered as an inadvertently ``treated'' group and compared with their compatriots residing in other areas of the state not part of that DMA. Limiting ourselves to these non-battleground states also allows us to decouple effects of advertising from the ground campaign and presidential visits \textemdash both of which are primarily geared towards battleground residents \textemdash that may influence citizens. The strategy yields what we call `No-intent to treat' Intent to Treat estimates of the effect of ads.

To estimate the impact of televised ads, we merged the NAES RCS data with the Wisconsin Ads Project data for 2000 and 2004 at the DMA level. The Wisconsin Ads Project data use satellite technology to keep record of all campaign ads broadcast. These ads are also content-coded for negativity, among other ad characteristics \citep[see][]{goldstein2007}. Using the ad data, we created a rolling thirty day window tallying all ads broadcast in the previous thirty days in each DMA. We then linked these data to respondent by date of interview and respondent's DMA of residence. We next implement the identification strategy, developed by \citet{huber2007}, to estimate the Average Treatment Effect of exposure to campaign advertising. To do this we estimated a hierarchical model predicting in-party minus out-party candidate trait and thermometer scores from as a function of total negative and all other advertising. In this specification we controlled for strength of partisanship,  which ought to be highly correlated with candidate evaluations, and a variety of demographic variables. We used fixed effects for states, so as to make only within state comparisons, and account for any state level unobserved variation, and random effects by DMA to account for clustering by DMA.

Letting $s$ indexing states, $j$ DMAs, and $i$ respondents, and letting $y$ denote affect towards candidate, $X$ as total ads shown, $\gamma$ the dependence between observations within a DMA, and $\delta$ the state fixed effects, formally, our model takes the following form:  
\begin{gather*}
  y_{ij}  \sim N(X_j\beta + \gamma_j + \delta_s; \sigma^2_{\epsilon}) \\
 \gamma_j  \sim N(0,\sigma^2_{\gamma})
\end{gather*}
\subsubsection*{Measures}
\textit{Battleground State Residence:} Our classification of battleground states is limited to those states where Election Day polling indicated the two candidates were within 5\% of each other. For 2008, the list of battleground states included CO, FL, IA, MI, MN, NC, NM, NH, NV, OH, OR, PA, WI, WV \citep[see][]{jackman2009}.  For 2004, the set included FL, IA, MO, MN, NH, NM, NV, OH, OR, PA, WI \citep[see][]{hillygus2005}. The 2000 list was identical to 2004 except for the addition of TN.

\textit{Candidate Favorability:} In all three NAES datasets, respondents were asked how favorably they rated each of the presidential candidates on a semantic scale that ran from ``very unfavorable'' to ``very favorable.'' In 2000, the scale had 101 points, while in 2004 and 2008 respondents were offered an 11 point scale. We rescaled the measures to range between 0 (``very unfavorable'') and 1 (``very favorable'') and created a difference score by subtracting the rating of the out party from in party candidate.

\textit{Candidate Traits:}  Respondents indicated the extent to which various traits described (2000) or applied (2004 and 2008) to the presidential candidates on a semantic scale that ranged from ``extremely well'' to ``not well at all'' (``does not apply at all'' in 2004 and 2008).  The scale consisted of four points in 2000, and ten in 2004 and 2008.  Our analysis is limited to those traits asked throughout the campaign season. In 2000, this set included ``really cares about people like me,'' ``honest,'' ``inspiring,'' and ``knowledgeable.''  In 2004, the common traits included ``cares about people like me,'' ``inspiring,'' ``strong leader,'' ``trustworthy,'' ``knowledgeable,'' and ``reckless.''  Finally, for 2008, we examine ``strong leader,'' ``trustworthy,'' ``has the experience needed to be president,'' and ``has the judgment needed to be president.''  

We rescored responses to positive traits to lie between 0 (``not well'') to 1 (``extremely well'') and reverse coded responses to the negative trait items.  We then computed a difference score by subtracting the out the party candidate trait rating from the in party candidate rating. The average correlation between the difference scores was at least .7 in each of the years, and consistently high over the duration of the campaign. On the thinking that such consistent ratings of one candidate over the other were reflective of affect, we chose to average the difference scores (see pg. 131, \citet{bartels2002} for a brief but more extensive reasoning behind why differences in trait ratings are unlikely to be caused by anything other than partisan perceptual bias, whose main antecedent variable is partisan affect).

\textit{Ideological Placement of Self and Candidates:} In all three NAES datasets included here respondents placed themselves and the candidates on a five-point ideology scale that ranged from ``very conservative'' to ``very liberal.'' We rescored the responses to the placement scale so that they ranged between 0 (``very conservative'') and 1 (``very liberal'').  We also took the absolute difference between respondents' self-placements and their placements of the in- and out party candidates.

\subsection*{Results}
For each of the three years under examination, favorability of the in-party candidate vis-\`{a}-vis the out-party candidate increased dramatically between the start of the campaign and Election Day (see Figure~\ref{fig:fig1}). The extent of this increase grew over the time period under examination, reaching a difference of more than 25 percentage points between the start of the campaign and Election Day in 2008. Tallying change in affect towards the in-party and the out-party candidate separately, we find that the affect towards out-party candidate declines at a faster clip than it rises for the in-party candidate ($p < .001$).

\begin{figure}[htpb!]
  \caption{Difference between In-Party and Out-Party Candidate Favorability Ratings Over the Campaign}
  \label{fig:fig1}
  \centering
    \includegraphics[width=1.0\textwidth]{../figs/inout_fav_naes_combo.pdf}
\end{figure}

While favorability ratings may be seen as capturing both ideological disagreement and ill-founded (if not ill-explained) aversive sentiment, differences in trait ratings are mostly an indicator of extent of partisan bias, unless as \citet{bartels2002} notes, ``Democrats and Republicans had markedly different views about what it means to be knowledgeable, and that each of the presidential candidates just happened to embody his own partisans' distinctive epistemological values, and that he did so increasingly as Election Day approached.'' So increases in differences in trait ratings between candidates are clearer indicators of affective mechanisms such as increased identity salience at work. As it turns out, a close facsimile of the pattern we saw for favorability ratings also holds for partisans' trait ratings of candidates. Difference between trait ratings of the in-party candidate and the out-party candidate increase sharply over the course of the campaign (see Figure~\ref{fig:fig4}). Again the sharpest rise was observed in 2008, and once again trait ratings of out-party candidate decline at a faster clip than they rise for the in-party candidate ($p < .001$).
\begin{figure}[htpb!]
  \caption{Difference between in-party and out-party candidate positive trait ratings over the campaign}
  \label{fig:fig2}
  \centering
    \includegraphics[width=1.0\textwidth]{../figs/inout_trt_naes_combo.pdf}
\end{figure}
To identify the role of exposure to the campaigns in this rise, we compared the trend in battleground and non-battleground states by regressing partisan affect on battleground, days to election, and the interaction between battleground residence and days to election. As expected, we found that the over-time differences in the in- and out-party candidate favorability ratings were greater in battleground states in 2004 and 2008. In the case of the 2000 campaign, however, the interaction between days before the election and battleground status proved non-significant. 

Next, we estimated the impact of televised advertising on trait ratings of candidates and favorability towards the candidates. Columns 1 and 3 of Table~\ref{tab:tab1} show estimates for difference in favorability between in- and out-party candidates, while columns 2 and 4 show estimates for difference in trait ratings of in- and out-party candidates. As we can see from Table~\ref{tab:tab1}, in both 2000 and 2004, exposure to negative advertising consistently polarizes affect towards presidential candidates, and evaluation of their traits. Individuals exposed to negative advertising come to hold larger gaps in how they feel towards their own-party's candidate and the out-party's candidate. The effect of all other political advertising is more mixed \textemdash it never has any influence on perceptions of differences in traits, and it increases the gap in affect in 2004, while reducing it in 2000.

\input{../tabs/ha_00_04_sg.tex}

One possible explanation for the increase in difference between favorability ratings of candidates over the course of the campaign is that voters learn about the candidates' positions on issues.  At the outset most voters have no idea where the candidates stand, but over time they are able to place the candidates on major issues.  And it is the greater awareness of the policy or ideological differences between the self and the candidates that causes people to express greater dislike for the out-party candidate.   We start our assessment of this learning-based explanation by comparing the mean ideological placements of the candidates over the course of the campaign (see Figures~\ref{fig:fig2} and \ref{fig:fig3}).  In-party placements for both Democrats and Republicans barely moved, while mean out-party candidate placements remained roughly constant for Democrats and showed only a slight trend among Republicans.
\begin{figure}[htpb!]
  \caption{In-party candidate placement over the campaign}
  \label{fig:fig3}
  \centering
    \includegraphics[width=1.0\textwidth]{../figs/out_rd_libcon_naes_combo.pdf}
\end{figure}

\begin{figure}[htpb!]
  \caption{Out-party candidate placement over the campaign}
  \label{fig:fig4}
  \centering
    \includegraphics[width=1.0\textwidth]{../figs/in_rd_libcon_naes_combo.pdf}
\end{figure}

The evidence presented in Figures~\ref{fig:fig3} and \ref{fig:fig4} shows that voters are not learning much about the ideology of their party's candidate.  In the case of the out-party candidate, any observed learning may be contaminated by projection, i.e. voters imputing greater ideological distance between themselves and the candidate they dislike.  The only significant trend that we do see is that the placements of Democratic candidates by Republicans become more extreme. This lone trend is consistent with affect-driven inference, when coupled with the lack of change in in-party ratings (which also consistently remain extremely close to respondent's self-placement on the ideology scale). Another explanation for the rise in affect over the course of campaign that we see may be strengthened partisan identities over the course of the campaign. As shown in ~\ref{si}, the proportion of partisans claiming a ``strong'' identity shows a small but significant gain over the course of the campaign, especially toward the closing stages.  The pattern is apparent for both Democrats and Republicans and is noticeably stronger in 2008 than in 2000.
\begin{figure}[htpb!]
  \caption{Absolute Difference Between self-placement and in and out party candidate placement on the ideology scale}
  \label{fig:fig5}
  \centering
    \includegraphics[width=1.0\textwidth]{../figs/out_rd_libcon_naes_combo.pdf}
\end{figure}
Of course, strengthening of partisan identities may have to do with learning of policy differences between parties and candidates over the course of the campaign.  But our aim here is not look into these potential mediators, but track how strength of partisan identities evolves over the campaign.  

\section*{Discussion}
There is strong evidence that over the course of campaigns, partisans' feelings towards the candidates and the parties, and their evaluations of the candidates become polarized. Negative messages are especially effective at causing this. The data also suggest that the rise in affect over the campaign is not primarily due to learning about real ideological positions of the candidates and the parties. The more likely explanation for the rise in partisan affect is strengthened partisan attachments, more salient partisan identities, and partisan processing of messages in a manner that persuade partisans about the negative traits of the out-party candidate, and positive traits of her own party. Our results add to the body of evidence that campaigns mobilize partisan identities and increase partisan animus. For instance, \citet{collins2011b} finds that exposure to negative political ads, particularly those that include partisan primes, mobilizes partisans to contribute to campaigns.

Our research also connects to ongoing research showing that partisanship is an increasingly salient basis for group identity \citep{iyengar2012, freeze2012}. The increased salience of partisanship as a basis for identity has come about, we suspect, because of the nature of modern political campaigns: they are increasingly long and loud. Campaigns have grown to occupy the entire period between election and re-election; campaigns are now permanent - each day devoted to partisan point scoring and ``winning the message'' \citep[see][]{blumenthal1982}.

Next, we discuss some limitations of our research. One concern about our evidence is generalizability to future campaigns. The campaign strategy of partisan reinforcement owes much to the nature of the current electorate. If we are sure about anything, it is that campaign strategies will evolve to keep up with the changing electorate. Then the question arises as to the prospect of change in American electorate. We cannot of course give a firm answer but we can point to the fact that on some crucial dimensions, as they relate to generalizability of our findings, American electorate has changed surprisingly little over the last fifty years - it remains largely disengaged with the issues, knows little, and tends to have strong partisan attachments. If past is to be taken as a predictor of future, we are in for much the same in terms of campaigns. 
    
Next concern relates not to the strategy but to dosage. There is little way we can calibrate the message ``dose'' in a campaign ad across times, and races. There are just too many variables: from affective reactions to background imagery and sound, to messages and specific context of each election and time during an election. It is likely that as a consequence of the data revolution, and burgeoning social scientific research, campaigns will learn to target individuals better, and will be able to deliver more ``potent'' messages to them. All in all, we think the effects we see here are likely to become stronger.

The above considerations deal with how campaigns are liable to evolve.  While our identifications strategy is relatively robust, avoiding the pitfalls of relying on self-reports of exposure, there is one reason to think that our estimates of the polarizing effects of televised campaigns are downwardly biased.  There is some evidence that the content, style, and volume, of television ads is context dependent and varies depending on the DMA they run in. For instance, variation in volume, content, etc. of ads within states (after accounting for say cost of running ads in the DMA) can be explained by the partisan profile of the DMA, among other things. If this is the case, the `spillover' design is going to underestimate (or merely mis-estimate) the impact of presidential ads because there is no relation between ads and partisan construction of the accidental spillover area. More generally, by decoupling strategic targeting (an essential part of campaigns) from message volume and content, we only get an estimate of the `treatment' (e.g. negative advertising) targeted crudely. If one wants an estimate of `strategic treatment' (which can be thought of as one kind of confounding variable), quasi-experimental designs relying on accidental spillover may be inappropriate.  For estimating the impact of strategically targeted televised campaign ads we first need to know how ad targeting depends on characteristics of the viewing audience - such as political interest \textemdash that moderate the impact of political ads \citep[see, for instance][]{ansolabehere1995}.  We can then estimate the effect of messages using the spillover design, but re-weight the effect using estimates of how the ad is targeted.  In future work, we hope to present such reweighted estimates.

\clearpage
\bibliographystyle{apsr}
\bibliography{polarcamp}  

\clearpage
\appendix
\renewcommand{\thesection}{SI \arabic{section}}
\renewcommand\thetable{\thesection.\arabic{table}}  
\renewcommand\thefigure{\thesection.\arabic{figure}}
\counterwithin{figure}{section}

\section{Strength of Partisan Identification Over The Campaign}
\label{si}

\begin{figure}[htpb!]
  \caption{Strength of Partisan Identification Over The 2000 Campaign}
  \label{fig:fig5}
  \centering
    \includegraphics[width=1.0\textwidth]{../figs/pidstr_2k.png}
\end{figure}

\begin{figure}[htpb!]
  \caption{Strength of Partisan Identification Over The 2004 Campaign}
  \label{fig:fig5}
  \centering
    \includegraphics[width=1.0\textwidth]{../figs/pidstr_2004.png}
\end{figure}

\begin{figure}[htpb!]
  \caption{Strength of Partisan Identification Over The 2008 Campaign}
  \label{fig:fig5}
  \centering
    \includegraphics[width=1.0\textwidth]{../figs/pidstr_2008.png}
\end{figure}

\end{document}
